% GENERATED_BY: oc_core_monograph_translation_draft_pipeline_v1.py
% AI_ASSISTED_TRANSLATION_DRAFT: TRUE
% THEOREM_REVIEW_STATUS: PENDING
% THEOREM_REVIEW_MODE: FORMAL_AI_GATE__CODEX_CLI_CHATGPT
% THEOREM_REVIEW_REVIEWER_ID:
% THEOREM_REVIEW_AUTHORITY_CLASS:
% THEOREM_REVIEWED_AT:
% THEOREM_REVIEW_DOSSIER_REF:
% SOURCE_LANGUAGE: EN
% TARGET_LANGUAGE: RU
% SOURCE_EN_REF: content/k_levels/k6.tex
% ================================================================
% ==== FILE: content/k_levels/k6.tex
% ================================================================

% ==============================
%  Ontology of Continua — Core
%  K-level Module: K6
%  FULL MODULE — FINAL
% ==============================

\subsubsection{\texorpdfstring{$K_6$}{K_6} Обзор}
\label{sec:k6-overview}

$K_6$ — это уровень, на котором \emph{когнитивная организация} возникает
как онтологическая размерность. Возбудимая электрическая динамика $K_5$
порождает устойчивые аттракторы; эти аттракторы становятся
\emph{моделями}, которые кодируют внутреннюю структуру, память,
предиктивные состояния и репрезентативные отношения.

Возникновение $K_6$ происходит, когда циклы возбуждения $K_5$ уже не
просто электрические волны, а становятся \emph{интерпретируемыми
информационно}, различия которых нельзя выразить в $A(K_5)$. Это
требует нового набора осей и новой размерностной структуры.

Минимальные признаки:
\begin{itemize}
    \item репрезентативные состояния,
    \item способность предсказания (ось ошибки предсказания),
    \item семантическая и связывающая структура,
    \item внутренние модели $\mu_t$,
    \item консолидация/реконсолидация памяти,
    \item устойчивые компоненты когнитивных графов,
    \item возникновение протосимволов.
\end{itemize}

$K_6$ — наименьший континуум, заслуживающий имя «когнитивный». Это ещё
не языковой или социальный уровень (они относятся к $K_7$).

% ================================================================
\subsubsection{Пространство состояний \texorpdfstring{$\Omega(K_6)$}{\Omega(K_6)}}

Состояние $K_6$:
\[
\Omega(K_6) =
\{
m,\; \mu_t,\; G_{\mathrm{cog}},\
x(t),\; r_i,\
PE_t,\; B_t,\
M_{\mathrm{con}}, M_{\mathrm{rec}}
\}.
\]

Элементы:
\begin{itemize}
    \item $m$ — активная когнитивная модель,
    \item $\mu_t$ — распределение состояний модели,
    \item $G_{\mathrm{cog}}$ — когнитивный граф (узлы = модели, рёбра = связывание),
    \item $x(t)$ — активационные паттерны, унаследованные от аттракторов $K_5$,
    \item $r_i$ — репрезентативные единицы (протосимволы),
    \item $PE_t$ — ошибка предсказания,
    \item $B_t$ — конфигурация связывания,
    \item $M_{\mathrm{con}}$, $M_{\mathrm{rec}}$ — сконсолидированные и
          реконсолидирующие состояния памяти.
\end{itemize}

Допустимость:
\[
D(m) \le \Theta_{\mathrm{cog}}, \qquad
|PE_t| < \Theta_{\mathrm{pred}}, \qquad
B_t \in \text{ёмкость связывания}.
\]

$D(m)$ — мера неконсистентности модели $m$.

% ================================================================
\subsubsection{\texorpdfstring{$\partial\Omega(K_6)$}{\partial\Omega(K_6)} (граница)}

Когнитивные границы включают:
\begin{itemize}
    \item ограничения репрезентативной ёмкости,
    \item границы глубины связывания,
    \item предельный горизонт предсказания,
    \item пределы устойчивости памяти,
    \item границы семантической когерентности,
    \item переходы, где $PE_t$ нельзя согласовать.
\end{itemize}

Переход через $\partial \Omega(K_6)$ приводит к коллапсу моделирования.

% ================================================================
\subsubsection{Оси \texorpdfstring{$A(K_6)$}{A(K_6)}}

Как установлено в the K6 construction route, минимальный набор:
\[
A(K_6) = \{
A_{\mathrm{sel}},\;
A_{\mathrm{cmp}},\;
A_{\mathrm{bind}},\;
A_{\mathrm{cat}},\;
A_{\mathrm{mem}},\;
A_{\mathrm{pred}},\;
A_{\mathrm{model}}
\}.
\]

Интерпретация:
\begin{itemize}
    \item $A_{\mathrm{sel}}$ — ось отбора (какая модель активна),
    \item $A_{\mathrm{cmp}}$ — ось сравнения,
    \item $A_{\mathrm{bind}}$ — ось связывания (композиция признаков),
    \item $A_{\mathrm{cat}}$ — ось категоризации,
    \item $A_{\mathrm{mem}}$ — ось памяти,
    \item $A_{\mathrm{pred}}$ — ось предсказания,
    \item $A_{\mathrm{model}}$ — ось идентичности модели.
\end{itemize}

Эти оси не сводятся к различиям возбуждения в $K_5$.

% ================================================================
\subsubsection{Потенциалы \texorpdfstring{$P(K_6)$}{P(K_6)}}

Потенциалы регулируют артефакты когниции:

\[
P(K_6)=
\{
P_{\mathrm{pred}},\;
P_{\mathrm{val}},\;
P_{\mathrm{bind}},\;
P_{\mathrm{sal}},\;
P_{\mathrm{mem}},\;
P_{\mathrm{elim}}
\}.
\]

Где:
\begin{itemize}
    \item $P_{\mathrm{pred}}$ — потенциал предсказания, управляющий
          обновлением модели,
    \item $P_{\mathrm{val}}$ — потенциал валентности (аффективное формирование),
    \item $P_{\mathrm{bind}}$ — потенциал связывания признаков,
    \item $P_{\mathrm{sal}}$ — модуляция салентности,
    \item $P_{\mathrm{mem}}$ — давление консолидации памяти,
    \item $P_{\mathrm{elim}}$ — устранение неконсистентных моделей.
\end{itemize}

% ================================================================
\subsubsection{Пороги \texorpdfstring{$\Theta(K_6)$}{\Theta(K_6)}}

Пороги определяют когнитивную жизнеспособность:
\[
\Theta(K_6)=
\{\Theta_{\mathrm{cog}},\;
\Theta_{\mathrm{pred}},\;
\Theta_{\mathrm{bind}},\;
\Theta_{\mathrm{sal}},\;
\Theta_{\mathrm{model}},\;
\Theta_{\mathrm{mem-stab}},\;
\Theta_{\mathrm{noise-cog}}\}.
\]

Значение:
\begin{itemize}
    \item $\Theta_{\mathrm{cog}}$: глобальная толерантность к
          неконсистентности,
    \item $\Theta_{\mathrm{pred}}$: максимально допустимая ошибка
          предсказания,
    \item $\Theta_{\mathrm{bind}}$: предел ёмкости связывания,
    \item $\Theta_{\mathrm{sal}}$: порог насыщения салентности,
    \item $\Theta_{\mathrm{model}}$: порог коллапса жизнеспособности модели,
    \item $\Theta_{\mathrm{mem-stab}}$: порог устойчивости памяти,
    \item $\Theta_{\mathrm{noise-cog}}$: толерантность к когнитивному шуму.
\end{itemize}

Переход этих порогов приводит к коллапсу $K_6$ или переходу к $K_7$ в
зависимости от направления.

% ================================================================
\subsubsection{Потоки \texorpdfstring{$J(K_6)$}{J(K_6)}}

Потоки информационные и репрезентативные:

\begin{itemize}
    \item $J_{\mathrm{pred}} = -\nabla_m PE_t$ — поток обновления
          предсказания,
    \item $J_{\mathrm{bind}}$ — переходы связывания/развязывания,
    \item $J_{\mathrm{cmp}}$ — динамика сравнения,
    \item $J_{\mathrm{sel}}$ — поток отбора модели,
    \item $J_{\mathrm{mem}}$ — поток консолидации/реконсолидации,
    \item $J_{\mathrm{elim}}$ — устранение неконсистентных моделей.
\end{itemize}

Стабильность требует:
\[
\sigma(J_{\mathrm{pred}}) < \sigma_{\max}(K_6).
\]

% ================================================================
\subsubsection{Циклы \texorpdfstring{$C(K_6)$}{C(K_6)}}

Характеристические циклы включают:

\[
C(K_6) =
\{
C_{\mathrm{pred}},\;
C_{\mathrm{bind}},\;
C_{\mathrm{mem}},\;
C_{\mathrm{cmp}},\;
C_{\mathrm{model}}
\}.
\]

Описания:
\begin{itemize}
    \item $C_{\mathrm{pred}}$: цикл предсказания (модель → предсказание →
          ошибка → обновление),
    \item $C_{\mathrm{bind}}$: цикл связывания признаков (композиция →
          стабилизация → освобождение),
    \item $C_{\mathrm{mem}}$: цикл памяти (консолидация → хранение →
          реконсолидация),
    \item $C_{\mathrm{cmp}}$: цикл сравнения–дифференциации,
    \item $C_{\mathrm{model}}$: цикл выбора, оценки и элиминации модели.
\end{itemize}

Каждый имеет характерный период $\tau_{C_j}$.

% ================================================================
\subsubsection{Время \texorpdfstring{$\tau(K_6)$}{\tau(K_6)}}

Время возникает из \emph{стабильности когнитивных циклов}:
\[
\tau(K_6)
= \min_j \tau_{C_j},
\qquad
\Pi(C_j) > \Theta_{\mathrm{time}}.
\]

Когнитивное время медленнее и более интегративно, чем электрическое
время ($K_5$).

% ================================================================
\subsubsection{Континуальность \texorpdfstring{$k(K_6)$}{k(K_6)}}

По общей формуле (the K6 construction route):

\[
k_6
= H(\Omega(K_6))
  \cdot \frac{|C_{\max}^{(6)}|}{|C_{\mathrm{all}}|}
  \cdot \frac{r}{r_{\max}}
  \cdot \left(1 - \frac{D(m)}{\Theta_{\mathrm{cog,max}}}\right)_+
  \cdot \left(1-\frac{\sigma(J)}{\sigma_{\max}}\right).
\]

Где:
\begin{itemize}
    \item $|C_{\max}^{(6)}|$ — размер крупнейшей связной компоненты в
          $G_{\mathrm{cog}}$,
    \item $r$ — число активных осей относительно максимального числа осей,
    \item $D(m)$ — метрика неконсистентности модели.
\end{itemize}

Интерпретация:
\begin{itemize}
    \item большая связная семантическая структура увеличивает $k_6$,
    \item неконсистентность и ошибка предсказания снижают $k_6$,
    \item фрагментация когнитивного графа коллапсирует $k_6$.
\end{itemize}

% ================================================================
\subsubsection{Структурное напряжение \texorpdfstring{$T(K_6)$}{T(K_6)}}

Источники:
\begin{itemize}
    \item высокая ошибка предсказания $PE_t$,
    \item несовместимое связывание,
    \item семантические конфликты,
    \item нестабильная реконсолидация памяти,
    \item некогерентная конкуренция моделей,
    \item разгонная салентность.
\end{itemize}

Коллапс при:
\[
T(K_6) > \Theta_{\mathrm{cog}}.
\]

% ================================================================
\subsubsection{Энергия \texorpdfstring{$E(K_6)$}{E(K_6)}}

Энергия является когнитивно-нейронным гибридом:
\[
E(K_6)
= E_{K_5}
+ E_{\mathrm{pred}}
+ E_{\mathrm{bind}}
+ E_{\mathrm{mem}}
+ E_{\mathrm{cmp}}
+ E_{\mathrm{model}}.
\]

Интерпретация:
\begin{itemize}
    \item энергия обновления предсказания,
    \item энергия поддержания памяти,
    \item затраты на связывание,
    \item энергия в репрезентативных графах,
    \item стоимость переключения модели.
\end{itemize}

% ================================================================
\subsubsection{Операторы на \texorpdfstring{$K_6$}{K_6} (\texorpdfstring{$\Psi$}{\Psi}, \texorpdfstring{$\Phi$}{\Phi}, \texorpdfstring{$\Lambda$}{\Lambda}, \texorpdfstring{$U$}{U}, \texorpdfstring{$\Chi$}{\Chi})}

\begin{itemize}
    \item $\Psi_{6\to7}$ — возникновение коммуникативных осей (общих моделей),
    \item $\Phi$ — эволюция репрезентативной структуры,
    \item $\Lambda$ — объединение компонентов модели в структуру более
          высокого порядка,
    \item $U$ — интеграция когнитивных циклов в устойчивые семантические
          единицы,
    \item $\Chi$ — ветвление в специализации (ориентация на память,
          предсказание или связывание).
\end{itemize}

% ================================================================
\subsubsection{Процессы на \texorpdfstring{$K_6$}{K_6}}

Типичные процессы:
\begin{itemize}
    \item формирование и исчезновение моделей,
    \item предсказание и коррекция ошибки,
    \item семантическое связывание,
    \item консолидация и реконсолидация,
    \item модуляция салентности,
    \item извлечение признаков,
    \item генерация протосимволов $\sigma$,
    \item аттракторная динамика, формирующая когнитивную топологию.
\end{itemize}

% ================================================================
\subsubsection{Прогнозы для \texorpdfstring{$K_6$}{K_6}}

\begin{enumerate}
    \item Когнитивный коллапс происходит, когда $PE_t$ накапливается быстрее,
          чем связывание может стабилизировать.
    \item Стабильная долговременная память требует циклов реконсолидации.
    \item У оси связывания есть максимальная ёмкость, которую можно
          экспериментально измерить.
    \item Семантическая связанность предсказывает устойчивость к шуму.
    \item Когнитивное время $\tau(K_6)$ растёт с ростом
          репрезентативной сложности.
    \item Прото-символьные состояния естественно появляются при достаточной
          связности $G_{\mathrm{cog}}$.
\end{enumerate}

% ================================================================
\subsubsection{Эксперименты для \texorpdfstring{$K_6$}{K_6}}

Возможные эмпирические аналоги:
\begin{itemize}
    \item модели нейронных сетей с аттракторами и контролируемым шумом,
    \item эксперименты предиктивного кодирования с измерением порогов $PE_t$,
    \item поведенческие тесты ёмкости связывания,
    \item протоколы восстановления реконсолидации памяти,
    \item тесты фрагментации семантического графа,
    \item наблюдение появления протосимволов в нейронных системах.
\end{itemize}

% ================================================================
\subsubsection{Коллапс и смерть \texorpdfstring{$K_6$}{K_6}}

Сбой, когда:
\[
\Omega(K_6) = \varnothing.
\]

Механизмы:
\begin{itemize}
    \item неконтролируемый рост ошибки предсказания,
    \item катастрофическая неудача связывания,
    \item коллапс семантического графа,
    \item распад памяти,
    \item каскады разгонной салентности,
    \item неспособность разрешить конфликты модели.
\end{itemize}

Смерть $K_6$ препятствует переходу к социальной когниции $K_7$.

% ================================================================
\subsubsection{Фальсифицируемость \texorpdfstring{$K_6$}{K_6}}

Модель предсказывает:
\begin{itemize}
    \item существование строгих порогов ошибки предсказания,
    \item необходимость лимитов ёмкости связывания,
    \item измеримые когнитивные циклы (предсказание, память, связывание),
    \item иерархическое возникновение протосимволов,
    \item зависимость семантической устойчивости от связности графа.
\end{itemize}

Опровержение возможно при наблюдении когнитивных систем без этих структур.

% ================================================================
\subsubsection{Разветвление / Онтологическое положение \texorpdfstring{$K_6$}{K_6}}

Ветвления:
\begin{itemize}
    \item системы, доминирующие памятью vs доминирующие предсказанием,
    \item высоко-связывающие vs низко-связывающие организмы,
    \item архитектуры богатыми vs бедными моделями,
    \item распределённые vs локализованные семантические графы.
\end{itemize}

Онтологическое позиционирование:
\[
K_5 \to K_6 \to K_7.
\]

$K_6$ — необходимое условие коммуникации, социального поведения и
общих моделей.

% ================================================================
\subsubsection{Связь с M-пространствами}

$M_6$ задаёт:
\begin{itemize}
    \item допустимые пейзажи ошибок предсказания,
    \item диапазоны топологии связывания,
    \item режимы устойчивости памяти,
    \item требования структурной когерентности,
    \item допустимую динамику репрезентативных графов.
\end{itemize}

Совместимость с $M_6$ определяет жизнеспособность когнитивных континуумов.
